% Generated draft: requires mathematical and source review.
\documentclass[11pt]{article}
\usepackage[margin=25mm]{geometry}
\usepackage{fontspec}
\setmainfont{Latin Modern Roman}
\usepackage{amsmath,amssymb,mathtools,unicode-math}
\setmathfont{Latin Modern Math}
\usepackage{enumitem,needspace,xurl,hyperref}
\setlength{\parindent}{0pt}
\setlength{\parskip}{5pt}
\setlength{\emergencystretch}{3em}
\newcommand{\N}{\mathbb{N}}
\newcommand{\Z}{\mathbb{Z}}
\newcommand{\Q}{\mathbb{Q}}
\newcommand{\R}{\mathbb{R}}
\newcommand{\C}{\mathbb{C}}
\newcommand{\F}{\mathbb{F}}
\newcommand{\A}{\mathbb{A}}
\newcommand{\PP}{\mathbb{P}}
\newcommand{\E}{\mathbb{E}}
\newcommand{\eps}{\varepsilon}
\newcommand{\dd}{\,\mathrm{d}}
\begin{document}
\section*{UnsolvedMath ID: 30006567}
\textbf{Source problem code:} catalogue entry
\subsection{Definitions and mathematical statement}
The intended objects are attached to a finite zero--one matrix \(A\). For a square instance \(A\in\{0,1\}^{n\times n}\), its classical automorphism group is
\[
\operatorname{Aut}(A)=\{\sigma\in S_n:P_\sigma A=AP_\sigma\},
\]
where \(P_\sigma\) is the permutation matrix of \(\sigma\). Regard this group as a classical compact quantum group. Write \(\operatorname{QAut}(A)\) for the graph quantum automorphism group intended in the cited report, and \(G_A^\infty\) for the report's limiting quantum isometry group associated with its Cuntz--Krieger construction and family \((G_A^\ell)_\ell\). The supplied excerpt does not give the defining relations of \(\operatorname{QAut}(A)\), the construction of \(G_A^\ell\), the meaning of the limiting operation, or the complete hypotheses on \(A\). These source-dependent objects therefore cannot be reconstructed unambiguously from this record alone.

For compact quantum groups \(H\) and \(G\), say that \(H\) is a quantum subgroup of \(G\) if there is a surjective unital \(*\)-homomorphism between their universal compact quantum group \(C^*\)-algebras,
\[
\pi:C^u(G)\longrightarrow C^u(H),
\qquad
(\pi\otimes\pi)\circ\Delta_G=\Delta_H\circ\pi,
\]
where \(\Delta_G\) and \(\Delta_H\) are the coproducts.

The intended question is: for every finite zero--one matrix \(A\) in the domain of the report's constructions, does
\[
\operatorname{QAut}(A)\ne\operatorname{Aut}(A)
\quad\Longrightarrow\quad
\operatorname{QAut}(A)\text{ is a quantum subgroup of }G_A^\infty
\]
hold? Here the inequality expresses the presence of quantum automorphisms beyond the classical automorphism group, using the report's identification of the classical group inside the quantum group. Equivalently, under that hypothesis, must there exist a surjective coproduct-preserving map
\[
C^u(G_A^\infty)\longrightarrow C^u(\operatorname{QAut}(A))?
\]
The qualification concerning the domain is necessary because the supplied clean statement says only ``finite zero--one matrix,'' whereas the omitted construction may require additional hypotheses. It is not a resolution of those missing hypotheses.

The excerpt supplies the first-level inclusion \(\operatorname{QAut}(A)\leq G_A^1\) as background and asks whether the corresponding inclusion holds for \(G_A^\infty\) in the nonclassical case. This draft records that question without asserting its answer.
\subsection{Short English statement}
Whenever the report's constructions are defined and \(A\) has quantum automorphisms beyond its classical automorphisms, must its full quantum automorphism group occur as a quantum subgroup of the limiting Cuntz--Krieger quantum isometry group \(G_A^\infty\)? The matrix hypotheses and the source-specific construction of the limiting group are missing from the supplied excerpt.
\subsection{Sources}
\begin{itemize}
\item Supplied local UnsolvedMath registry record, permanent ID \(30006567\), source problem code \texttt{catalogue entry}, entitled ``Quantum Automorphisms in Limiting Cuntz--Krieger Isometry Groups.'' The project identifies \texttt{data/unsolvedmath/problems.json} and \texttt{catalog/problems.json} as registry files; their contents were not independently inspected for this draft. No individual proposer is identified in the supplied record.
\item \emph{Noncommutative Harmonic Analysis and Quantum Information} (2026), Oberwolfach Reports, DOI \texttt{10.4171/owr/2026/4}, pp.~217--275. Supplied source and literature link: \url{https://doi.org/10.4171/owr/2026/4}. The relevant excerpt labels the subgroup question ``Question 1.''
\item Supplied literature assessment, checked on \texttt{2026-08-22}: the record labels the problem open and reports that the primary report explicitly describes the limiting subgroup inclusion as unknown when \(\operatorname{QAut}(A)\ne\operatorname{Aut}(A)\). This is supplied dated curation, not an independently verified literature assessment. The record also labels the statement \texttt{corrected\_verified} and says it was reconstructed and checked against the official report; that verification is supplied provenance, not verification performed here.
\item The background excerpt additionally cites [Ban05] concerning graph quantum automorphism groups, [GJ18] concerning quantum isometry groups of classical smooth manifolds, and Alexander Mang's thesis [Man22] concerning unitary easy quantum groups. Complete bibliographic details for these reference labels are not supplied.
\end{itemize}
\end{document}
